% !PRINTED 443
\weberpage{443}
Die rechte Seite dieser Formel ist aber von $m'$ ganz unabhängig.
Der Werth der linken Seite ist also derselbe, als ob
$m' = 1$ wäre, und wenn man für die Producte die Werthe aus
(4) setzt, so folgt

\[
\left(\frac{mm'}{n}\right) : \left(\frac{m'}{n}\right) = \left(\frac{m}{n}\right) : \left(\frac{1}{n}\right),
\]

oder

\[
8.\quad \left(\frac{mm'}{n}\right) = \left(\frac{m}{n}\right)\left(\frac{m'}{n}\right).
\]

Hier war zunächst vorausgesetzt, dass $m$ und $m'$ positiv und
ungerade seien. Nach 3. und 5. aber bleibt 8. auch noch richtig,
wenn $m$ oder $m'$ negativ oder gerade ist, wenn nur $m$ und $m'$
relativ prim zu $n$ sind.

Nehmen wir wieder $m$ ungerade und positiv an, so folgt
jetzt aus (4) und (9)

\weq{10}{\left(\frac{m}{n}\right) = (-1)^{\frac{(m-1)(n-1)}{4}} 2^{\frac{(m-1)(n-1)}{2}} \overset{\alpha}{\Pi}\overset{\beta}{\Pi} (\alpha - \beta).}

Wenn wir hierin $m$ mit $n$ vertauschen, so ändert in dem
Doppelproduct auf der rechten Seite jeder Factor sein Vorzeichen;
alles Andere bleibt ungeändert. Die Anzahl dieser Factoren
ist aber

\[
\frac{n-1}{2} \cdot \frac{m-1}{2}
\]

und daraus folgt

\[
9.\quad \left(\frac{m}{n}\right) = (-1)^{\frac{(m-1)(n-1)}{4}} \left(\frac{n}{m}\right),
\]

\gesperrt{oder} $\left(\dfrac{m}{n}\right)$ \gesperrt{ist gleich} $\left(\dfrac{n}{m}\right)$\gesperrt{, wenn von den beiden
ungeraden Zahlen} $m, n$ \gesperrt{wenigstens eine von der Form}
$4k+1$ \gesperrt{ist, und} $\left(\dfrac{m}{n}\right)$ \gesperrt{ist entgegengesetzt
 zu} $\left(\dfrac{n}{m}\right)$\gesperrt{, wenn beide Zahlen von der
Form} $4k+3$ \gesperrt{sind.}

Dieser berühmte Satz ist unter dem Namen des \gesperrt{Reciprocitätsgesetzes}
 bekannt. Der hier gegebene Beweis rührt von
\gesperrt{Eisenstein} her. Man kann mit seiner Hülfe und nach 5. und
6. den Werth des Symbols $\left(\dfrac{m}{n}\right)$ sehr schnell ermitteln, indem
%CONT>>
